\section{Linear-time phase extraction and full normalization}
\label{sec:phase}
This section gives complete written proofs, including the tape schedules;
its results are not formalized in Lean. A phase identifies a column's
position within its cyclic orbit. Computing that index and moving the
column to the indicated position are different output tasks. We prove
that the first is linear-time for a specific recurrence, and that the
second is equivalent to $\WORD$.

\subsection{The phase recurrence}
Throughout this section, bits within a word are ordered least significant
first. For a column of dyadic height $n$ and word width $w$, write
\[
 f(X)=\sum_{r=0}^{n-1}x_rX^r,\qquad x_r\in\mathbb F_2^w.
\]
Coefficients are vectors: the only coefficient operations used are XOR
and bit inspection. The cyclic shift $J_n$ moves source row $r$ to row
$r+1\pmod n$, so $(J_nf)_r=x_{r-1\bmod n}$. It is multiplication by
$X$ modulo $X^n-1$. Since $n$ is a power of two, substitution $X=1+y$
identifies this module with vector polynomials modulo $y^n$.

For $f\ne0$, write
\[
 f(1+y)=y^v(A+By+\cdots),\qquad A\ne0,
\]
where $v$ is the least nonzero coefficient degree and $B=0$ if $v=n-1$.
Let $p(f)$ be the bit of $B$ at the least significant set-bit position of
$A$. Set $p(0)=0$. A nonzero fixed column has $v=n-1$, and hence $p=0$.
For $n=2m$, define adjacent-row maps
\[
 (C_0f)_k=x_{2k}\mathbin\oplus x_{2k+1},\qquad
 (C_1f)_k=x_{2k+1}\quad(0\le k<m).
\]
Put $N(f)=J_n^{-p(f)}f$ and
$K(f)=C_{v\bmod2}N(f)$ for nonzero $f$; set $K(0)=0$.
The phase convention is
\begin{equation}\label{eq:phase-recurrence}
 \phi_n(f)=
 \begin{cases}
  0,&J_nf=f,\\
  p(f)+2\phi_m(K(f)),&J_nf\ne f.
 \end{cases}
\end{equation}
Every height-one column is fixed, so this defines $\phi$ at all dyadic
heights. It does not specify a lexicographically least rotation.

\begin{theorem}[Paid phase extraction]\label{thm:paid-phase}
\notformalized\ Let $d=2^h$, $h\ge0$, $b=\max(1,h)$, and $P=d^2b$.
There is one uniform deterministic native machine that, on every raw
row-major $d\times d$ matrix of $b$-bit entries, outputs in column order
the phase $\phi_d(f)$ and least positive cyclic period $L(f)$ of each
column. Its time is $O(P)$, including raw size discovery, all workspace
operations, and output. For every column,
\[
 0\le\phi_d(f)<L(f),\qquad
 \phi_d(J_df)=(\phi_d(f)+1)\bmod L(f).
\]
Each output record consists of an $h$-bit phase followed by an $(h+1)$-bit
period, both least significant first, without separators. Zero and fixed
columns have phase zero and period one.
\end{theorem}

\begin{proof}
\emph{Period and carry.}
Put $g(y)=f(1+y)$. Multiplication by $J_n$ becomes multiplication by
$1+y$ modulo $y^n$. For $f\ne0$, this preserves $v,A$. It fixes $f$
exactly when $v=n-1$; otherwise it replaces $B$ by $B+A$, so
$p(J_nf)=1-p(f)$. The least period divides $n$ and is therefore a power
of two. For every $2^s\le n$,
\[
 (J_n^{2^s}-I)g=y^{2^s}g\pmod {y^n}.
\]
For nonzero $g$, this vanishes exactly when $v+2^s\ge n$. Consequently
\begin{equation}\label{eq:phase-period}
 L(f)=\min\{2^s:2^s\ge n-v\}\quad(f\ne0),\qquad L(0)=1.
\end{equation}
For a nonfixed column the one-bit normalization satisfies
\[
 N(J_nf)=
 \begin{cases}N(f),&p(f)=0,\\J_n^2N(f),&p(f)=1.\end{cases}
\]

\emph{Shrinking child and phase induction.}
Separating even and odd raw coefficients gives
\begin{equation}\label{eq:phase-components}
 f(1+y)=(C_0f)(1+y^2)+y(C_1f)(1+y^2).
\end{equation}
Multiplication by $J_n^2$ multiplies both components by $1+y^2$ in
their respective quotients. Thus
$C_iJ_n^2=J_mC_i$ for $i=0,1$, including cyclic wrap.
The component selected by $v\bmod2$ has valuation $\floor{v/2}$ and
leading word $A$. For $L(f)>1$, \cref{eq:phase-period} therefore gives
\[
 L(K(f))=L(f)/2,
\]
because $m-\floor{v/2}=\ceil{(n-v)/2}$. For zero or fixed columns,
the child remains zero or fixed. The selected parity is invariant under
$J_n$, so for nonfixed columns
\[
 K(J_nf)=K(f)\quad\text{if }p(f)=0,\qquad
 K(J_nf)=J_mK(f)\quad\text{if }p(f)=1.
\]
Induction in \cref{eq:phase-recurrence} proves the range and equivariance.
In the second case, with $L'=L(f)/2$, the new phase is
$2((\phi_m(K(f))+1)\bmod L')$, which is the old phase plus one modulo
$2L'$. This includes $L'=1$. The fixed cases hold modulo one.

\emph{Extracting two coefficients without expanding the Taylor polynomial.}
The implementation needs only $v,A,B$ at each stage. For a nonzero
polynomial $f$ of degree less than $n=2m$, keep a boundary word $e$ and
an offset $t$. Split $f=a+X^mc$, with both halves of degree less than
$m$, and put $\ell=a+c$. Then
\[
 f(1+y)+ey^{2m}
 =\ell(1+y)+y^mc(1+y)+ey^{2m}.
\]
If $\ell\ne0$, replace $(f,e)$ by $(\ell,c(1))$ and leave $t$
unchanged. The leading degree is below $m$; the next coefficient crosses
the boundary only if that degree is $m-1$, in which case the needed
additional word is exactly $c(1)$. If $\ell=0$, replace $f$ by $c$,
retain $e$, and add $m$ to $t$, using the exact factorization
$y^m(c(1+y)+ey^m)$. Starting with $e=t=0$, these updates preserve
the first two leading words and preserve the original valuation as
$t$ plus the current valuation. At height one, the remaining word is
$A$, the boundary word is $B$, and $t=v$. A separate zero-column flag
keeps all fields zero for zero columns. The offset additions use distinct
descending powers of two and cannot overflow the original $h$-bit field.

\emph{Native implementation of the selector.}
Assume $h\ge1$, so $b=h$. A rectangular current bank has $n$ rows,
$d$ columns, and the original $b$-bit words. Copy it before selection,
so the current payload remains available. At each inner fold, sequentially
copy its first and last halves to two banks, XOR them to a third, and
accumulate, for each column, a nonzero flag for $\ell$ and the word
$c(1)$. Store these in a scratch row scanned alongside every data row
and rewound at each row end. Scan the metadata once to update $e,t$.
Then scan the $\ell$ and $c$ banks together, selecting the indicated
entry in each column; both source heads advance regardless of the choice.
Erase the used temporary spans and repeat with half as many rows.

A fixed pool of bank and metadata tapes suffices. A bank has an origin
marker and entry/row delimiters, adding only $O(1)$ cells per entry or
row. Entry rulers and counters occupy ordinary tapes. Their $O(h)$ work
per entry is $O(b)$; each scratch-row return costs $O(db)$ and is charged
to that data row. Separate metadata scans cost $O(db)$ per fold. With
$q=\log_2 n$, all selector copies, scans, erasures, and returns therefore
cost
\begin{equation}\label{eq:phase-selector-cost}
 O(ndb+db(q+1)).
\end{equation}
The fixed-column test uses the current height, $v=n-1$.

\emph{Native outer update.}
At each height $n=2m$, run this selector on a copy of the current bank.
Store the per-column bits $p$ and $v\bmod2$, using zero for the latter
on zero columns. Form a uniformly shifted bank $J_n^{-1}f$ by saving
the first row, copying the remaining rows, and appending the saved row.
Merge shifted and unshifted banks using the $p$ scratch row. This
implements $N(f)$ with a zero-or-one-row shift per column in $O(ndb)$
time; it does not apply an arbitrary phase for free.

Split the resulting bank stably into even-row and odd-row banks $E,O$.
Scanning both in row order, write $E\oplus O$ where $v$ is even and
$O$ where it is odd. The resulting row-major bank is exactly $K(f)$
column by column. Scratch scans and returns cost $O(ndb)$. Erase the
old bank and temporary spans before reusing their tapes, retaining only
this one child and the metadata. The number of tapes is independent of
$h$; both the inner and outer recurrences are loops over fixed tape pools.

Keep $d$ permanent output records. At outer level $i$, write $p$ into
phase bit $i$, by scanning each whole record if necessary; this costs
$O(dh)$ per level. Continue all columns for all $h$ levels: fixed and
zero children remain fixed and contribute zero bits. Use the first
selector's valuation to compute the \emph{original} period from
\cref{eq:phase-period}, and never overwrite it with a child period.
Even an $O(h^2)$ local binary routine per column suffices for this step.

\emph{Setup, total cost, and output.}
Count raw input length using a binary counter with its low end fixed.
An increment and return across $t$ trailing ones costs $O(t+1)$, giving
$O(P)$ total. Rewind the input. Test $P=h4^h$ for $h$ up to the bit
length of $P$ using ordinary binary tape arithmetic. This one-time setup
is polynomial in $\log(P+1)$ and hence $O(P)$ with one uniform constant.
Create the rulers, $O(h)$-bit counters, bank delimiters, and output records
by paid scans. No dimension is stored in finite control or supplied as advice.

Across outer heights $d,d/2,\ldots,2$, the data-bank cost is
$O(db\sum_i d/2^i)=O(P)$. The nested selector metadata contribute
\[
 O\left(db\sum_{q=1}^h(q+1)\right)=O(dbh^2)=O(P),
\]
using $h^2\le2\cdot2^h=2d$ for $h\ge1$.
Period computation and phase-record writes cost $O(dh^2)$ in total.
All bank erasures and scratch-row returns are included in these bounds.
Finally scan the records in column order and append exactly $d(2h+1)$
bits, with no workspace use of the output tape. The case $d=1$ outputs
the one-bit period $1$ and an empty phase field. Invalid lengths may halt
with a fixed fallback. This proves the promised uniform native bound.
\end{proof}

The theorem removes phase discovery as a possible source of superlinear
cost for this recurrence. It does not recover the original column from
its phase and period, nor does it output a rotated payload.

\subsection{The full normalizer is equivalent to word transposition}
Define the full column normalizer by
\[
 \mathcal N_{\phi,n}(f)=J_n^{-\phi_n(f)}f.
\]
Equivariance makes this constant on every cyclic orbit. The bank task
applies it to every column and writes the entire resulting raw row-major
matrix, with the original word width and no additional encoding.

\begin{lemma}[Exact marker phase]\label{lem:phase-marker}
\notformalized\ Let $n$ be a power of two. In a column of any word
width $w\ge1$, suppose the least significant bit is one exactly at
row $a$, $0\le a<n$, and zero elsewhere. All higher bits may be
arbitrary. Then $L(f)=n$ and $\phi_n(f)=a$.
\end{lemma}
\begin{proof}
The single marker forces period $n$. For $n\ge2$, its constant Taylor
coefficient is one, so $v=0$ and the pivot is the least significant bit,
regardless of the payload in higher bits. Its linear Taylor coefficient
is $a\bmod2$, hence $p=a\bmod2$. The shift $J_n^{-p}$ moves the
marker to the even row $a-p$; the selected $C_0$ child has its marker
at $\floor{a/2}$. Induction, with the immediate base case $n=1$, gives
$\phi_n(f)=(a\bmod2)+2\floor{a/2}=a$.
\end{proof}
The exact origin matters: equivariance alone would allow an unknown
orbit-dependent additive constant. This marker fixes that constant even
when all other bit lanes carry arbitrary data.

\begin{theorem}[Normalizer equivalence]\label{thm:normalizer-word}
\notformalized\ One uniform deterministic native machine computes
$\mathcal N_{\phi,d}$ on all raw dyadic $d\times d$ matrices with
$d=2^h\ge2$ and canonical width $b=h$ in $O(d^2b)$ time if and
only if $\WORD$ holds. Here $\phi$ is exactly
\cref{eq:phase-recurrence}, with its least significant set-bit pivot.
\end{theorem}
\begin{proof}
\emph{Native interfaces.}
Every intermediate bank below is a work tape. A fixed subroutine's
read-only input head and append-only output can be simulated on such
banks with constant transition overhead. End markers delimit the raw
input; an origin track and a side flag permit returning its input head
to origin after halting. The return costs at most its runtime plus the
bank length. Reserve separate initially blank workspace sets for the
constant number of calls. This uses a fixed number of tapes and does
not assume free cleanup or free resets. Only the wrapper's final answer
is written on the actual append-only output tape.

Raw size recovery costs $O(P)$ as in \cref{thm:paid-phase}.
For arbitrary canonical dimensions, binary search the strictly increasing
function $d\mapsto d^2\max(1,\ceil{\log_2d})$ after counting the input
length. Its one-time arithmetic uses polynomial time in $\log(P+1)$.
Ordinary $O(b)$-bit counters and $b$-cell rulers cost $O(b)$ per
entry visited. Copies, seeks, rewinds, and zero padding are charged by
the number of cells traversed or written.

\emph{From a normalizer to a column shear.}
First take $d=2^h\ge2$, $b=h$, and put $D=2d$. Given any matrix
$A$, construct a raw $D\times D$ matrix $B$ of canonical width $h+1$.
For each active column $c<d$, put a single least-significant-bit marker
at row
\[
 a_c=2((-c)\bmod d).
\]
For dummy columns $d\le c<D$, put that marker at row zero. In the
higher $h$ bits of $B[2r,c]$, for $r,c<d$, copy $A[r,c]$; all
other payload bits are zero. This is a sequential construction: at row
$2r$, compare $c$ to $(-r)\bmod d$, write the marker first, then
copy the payload. Write all dummy entries and all odd rows explicitly.
Every comparison costs $O(h+1)$ per entry, and
\[
 D^2(h+1)=4d^2(h+1)\le8P.
\]
Thus preparation and rewind cost $O(P)$.

By \cref{lem:phase-marker}, normalization gives
$C[s,c]=B[s+a_c\bmod D,c]$ for active columns. Scan its entire raw
output, keeping the first $d$ columns of the even rows and dropping
the first bit of each retained entry. The extracted bank satisfies
\[
 U[r,c]=A[r-c\bmod d,c].
\]
It is already in row-major order. This implements the source-motion
shear $V(r,c)=(r+c,c)$ with one normalizer call and $O(P)$ other work.

\emph{Completing the transpose and extending to all sizes.}
The row maps $H_-(r,c)=(r,c-r)$ and $R(r,c)=(r,-c)$ are linear-time
native operations. For $H_-$, copy each row to a reusable buffer and
emit its suffix starting at entry $r$, then its prefix. For $R$, emit
entries in the order $0,d-1,d-2,\ldots,1$, each word forward; moving
back two word lengths after an emission reaches the preceding entry.
Including seeks, returns, and counter updates, each routine costs
$O(db)$ per row. Tracking source positions through the chronological
sequence $H_-,V,H_-,R$ gives
\[
 (r,c)\longmapsto(r,c-r)\longmapsto(c,c-r)
       \longmapsto(c,-r)\longmapsto(c,r).
\]
It therefore computes dyadic transpose.

For arbitrary $d\ge2$, put $b=\ceil{\log_2d}$ and $N=2^b$.
Pad to the top-left corner of an $N\times N$ matrix with the same
word width, run the complete dyadic transpose, and crop its top-left
$d\times d$ corner. Since $d\le N<2d$, all padding and cropping
cost $O(P(d))$. The one normalizer call has raw length
\[
 4N^2(b+1)\le8N^2b<32P(d).
\]
Cropping occurs after the full transpose, not after an intermediate
shear. The case $d=1$ is a fixed copy routine. Constant-overhead
simulation and all wrappers thus prove $\WORD$.

\emph{From word transposition to the normalizer.}
On a dyadic bank $X$, run \cref{thm:paid-phase}, capturing its output
on a work tape and retaining $X$. Pay the input-head reset and rewind
the phase records. Transpose $X$ to a fresh bank $Y=X^{\mathsf T}$.
Row $c$ of $Y$ is original column $c$, and the phase records are in
this same sequential order. Copy that row to a reusable buffer, read
its phase $t$, and emit the suffix beginning at entry $t$, then the
prefix. Ignore its period field. This produces the row with entry $r$
equal to the old entry $r+t\bmod d$, in $O(db)$ time per row, including
every seek, counter, rewind, and buffer reuse. A second transpose gives
the raw bank whose column $c$ is $J_d^{-\phi_d(X_{:,c})}X_{:,c}$.
Zero and short-period columns need no exception beyond their computed
phases. There are two fixed transposer calls and $O(P)$ other work;
if their runtime is $t(P)=O(P)$, the total is $O(t(P)+P)=O(P)$.
\end{proof}

This equivalence supplies a concrete alternative construction target and
an equally strong impossibility target. It does not prove a transposer,
does not establish $\WORD\Rightarrow\BIT$, and does not assert the
same equivalence for an arbitrary choice of orbit representative.

\subsection{Why the discarded sibling still matters}
Phase extraction retains just one shrinking child. Full normalization
must retain enough information for both components. At height $n=2m$,
write $U=C_0N(f)$, $V=C_1N(f)$, and $q=\phi_m(K(f))$.
For fixed columns as well as nonfixed ones, the recurrence gives
$\phi_n(f)=p(f)+2q$. Commutation with $J_n^2$ yields
\[
 C_0\mathcal N_{\phi,n}(f)=J_m^{-q}U,\qquad
 C_1\mathcal N_{\phi,n}(f)=J_m^{-q}V.
\]
Both siblings require the \emph{same} phase $q$, selected by one of
them. Inverting the adjacent-row maps reconstructs the normalized even
entries by XOR of these two shifted siblings and the odd entries from
the shifted $V$. Independently normalizing each sibling need not work.

For example, with two-bit words and $f=(2,2,1,0)$, the recurrence gives
$v=0$, $p=0$, $U=(0,1)$, $V=(2,0)$, and $q=\phi_2(U)=1$.
The correct normalized column is $(1,0,2,2)$. But $\phi_2(V)=0$;
using each sibling's own phase reconstructs $(3,2,0,0)$ instead.
More generally, placing a single lowest-bit marker in $U$ and zeros
in that lane of $V$ leaves all higher lanes of both siblings independent.
There is no general recovery of the discarded payload from the selected
child and the phase alone.

Literal recursion on both full siblings preserves the total payload
volume at each level, so a full scan per level costs $O(P\log d)$.
That observation is not a lower bound for other computations. The
remaining constructive task is to share work across these coupled
updates while preserving and eventually emitting the entire raw payload.
The normalizer equivalence shows precisely why achieving this in linear
time would settle canonical word transposition.
